\documentclass[11pt,a4paper]{article}

\usepackage[margin=2.4cm]{geometry}
\usepackage{amsmath,amssymb,amsthm}
\usepackage{booktabs}
\usepackage{array}
\usepackage{hyperref}
\usepackage[expansion=false]{microtype}
\usepackage{enumitem}

\hypersetup{
  colorlinks=true,
  linkcolor=black,
  citecolor=black,
  urlcolor=blue,
  pdftitle={Deterministic Conservative Integer Routing on Discrete Lattices},
  pdfauthor={Sanket Hazra}
}

\newtheorem{theorem}{Theorem}
\newtheorem{lemma}{Lemma}
\newtheorem{proposition}{Proposition}
\newtheorem{corollary}{Corollary}
\theoremstyle{definition}
\newtheorem{definition}{Definition}
\newtheorem{remark}{Remark}

\newcommand{\Z}{\mathbb{Z}}
\newcommand{\N}{\mathbb{N}}
\newcommand{\R}{\mathcal{R}}
\newcommand{\Om}{\Omega}

\title{\bfseries Deterministic Conservative Integer Routing\\
on Discrete Lattices:\\[0.3em]
\large A Verified Reference Implementation with a Multilevel Extension}

\author{Sanket Hazra\\
\small Independent Researcher, Midnapore, West Bengal, India\\
\small \texttt{sanket8hazra@gmail.com}}

\date{Working paper --- 20 November 2026}

\begin{document}

\maketitle

\begin{abstract}
\noindent
This paper studies a single primitive: the deterministic distribution of a
non-negative integer among the $k$ ordered neighbours of a lattice cell such
that the quantity is preserved exactly and the allocation is as even as
integer arithmetic allows. The rule is $n = kq + r$ with $q$ given to every
destination and one additional unit to the first $r$ of them. We establish
exact conservation, a one-unit bound on allocation imbalance, determinism,
component bounds and translation invariance; we extend the result from a
single neighbourhood to a whole lattice under both sequential and simultaneous
update semantics; we show that conservation survives coincident neighbours on
a small periodic domain once the conservation scope is de-duplicated; and we
give a two-level restriction--prolongation scheme whose cycles preserve the
global total exactly. A reference implementation is described and verified by
88 automated tests covering 45{,}009 operator cases and full sweeps of
lattices up to $10^4$ cells at three neighbourhood spacings. We are explicit
about what is not established: no performance comparison against any existing
architecture, no hardware realisation, no physical interpretation, and no
claim of novelty for the routing rule itself, which is the largest-remainder
rule long known in apportionment theory. The contribution offered here is the
verified construction, the lattice-level conservation results, the multilevel
extension, and a characterisation of the rule family identified as the next
step.

\medskip
\noindent
\textbf{Keywords:} discrete lattice; integer routing; conservation;
quotient--remainder decomposition; multilevel transfer; verified reference
implementation.
\end{abstract}

\section{Introduction}

Numerical work on regular lattices is usually expressed in floating point, and
the resulting state is conserved only to within accumulated rounding. For some
problems that is the right trade: the physical model itself is continuous and
the discretisation error dominates. For others the quantity being moved is
genuinely discrete --- units allocated to processors, tokens assigned to
replicas, messages counted per route --- and an approximate conservation law is
a modelling error rather than a numerical one.

This paper takes the second case and asks a narrow question. Given a lattice
cell holding a non-negative integer, and a fixed ordered set of neighbours, is
there a routing rule that is (i) exactly conservative, (ii) as balanced as
integer arithmetic permits, (iii) deterministic, and (iv) invariant under a
uniform shift of the routed quantity? The answer is yes, the rule is
elementary, and the interesting content lies not in the rule but in what
happens when it is applied repeatedly across a lattice, at a boundary, and
across resolution levels.

The rule is not new. Distributing a remainder one unit at a time to the first
$r$ destinations is the largest-remainder method of apportionment, in use since
the eighteenth century and analysed at length by Balinski and Young
\cite{balinski1982}. We say so plainly at the outset because the opposite
practice --- presenting a known construction as a discovery --- is the most
common way work of this kind loses credibility. What we offer instead is a
verified implementation, a set of lattice-level conservation results that are
not immediate from the per-cell identity, a multilevel extension with an exact
conservation guarantee, and a precise statement of the open question that would
turn the description into a theorem (\S\ref{sec:future}).

\paragraph{What this paper does not claim.}
It makes no performance claim of any kind. It presents no hardware. It does not
interpret the lattice as spacetime or as a physical field, and it does not
claim that integer arithmetic is faster, smaller or more efficient than any
alternative. Section~\ref{sec:limits} states each of these explicitly, and
Appendix~\ref{app:withdrawn} records the claims from earlier versions of this
work that are withdrawn, together with the reason.

\paragraph{Structure.}
Section~\ref{sec:model} fixes notation and defines the operator.
Section~\ref{sec:props} proves the per-cell properties.
Section~\ref{sec:lattice} lifts them to a whole lattice under two update
semantics. Section~\ref{sec:boundary} treats boundaries and the coincident
neighbour case. Section~\ref{sec:multi} gives the multilevel scheme.
Section~\ref{sec:impl} describes the implementation and its verification.
Sections~\ref{sec:related}--\ref{sec:conclusion} place the work, state its
limits and identify what would come next.

\section{Model}\label{sec:model}

\begin{definition}[Lattice state]
Let $d \geq 1$ and $V \subseteq \Z^d$. A \emph{state} assigns to each
$v \in V$ a non-negative integer $\Om(v) \in \N$. The state is finitely
supported: $\Om(v) = 0$ for all but finitely many $v$.
\end{definition}

\begin{definition}[Axial neighbourhood]
For $v \in \Z^d$ and a spacing $s \in \Z_{\geq 1}$, the axial neighbourhood is
the ordered tuple
\begin{equation}
N_s(v) = \bigl(v + s e_1,\; v - s e_1,\; v + s e_2,\; v - s e_2,\;
\ldots,\; v + s e_d,\; v - s e_d\bigr),
\end{equation}
of length $k = 2d$, where $e_i$ is the $i$-th canonical basis vector. Every
element lies at $L_1$ distance $s$ from $v$. The ordering is part of the
definition: the rule below assigns the remainder positionally, so the
neighbourhood is an ordered tuple and not a set.
\end{definition}

For $d = 4$ we obtain $k = 8$; for $d = 2$, $k = 4$. Nothing in the analysis
depends on $d = 4$, and we state results for general $k = 2d$.

\begin{definition}[Routing operator]\label{def:op}
For $n \in \N$ and $k \geq 1$, write the quotient--remainder decomposition
\begin{equation}
n = kq + r, \qquad q = \left\lfloor \frac{n}{k} \right\rfloor, \qquad
0 \leq r < k .
\end{equation}
The routing vector $\R(n) = (R_1(n), \ldots, R_k(n))$ is
\begin{equation}\label{eq:rule}
R_i(n) = q + \mathbf{1}[\,i \leq r\,], \qquad i = 1, \ldots, k ,
\end{equation}
equivalently $r$ components equal to $q+1$ followed by $k-r$ components equal
to $q$. For $n = 27$ and $k = 8$,
\begin{equation}
\R(27) = (4,4,4,3,3,3,3,3).
\end{equation}
\end{definition}

\begin{definition}[Source transition]
Given a centre $v_c$ with $\Om_c$ and neighbours $N_1, \ldots, N_k$ with states
$\Om_i$, and an admissible quantity $n \leq \Om_c$,
\begin{equation}\label{eq:transition}
\Om_c' = \Om_c - n, \qquad \Om_i' = \Om_i + R_i(n) \quad (i = 1,\ldots,k).
\end{equation}
\end{definition}

The operator is a redistribution and not a replication: the source gives up
exactly the quantity that the destinations receive. Equation
\eqref{eq:transition} is the entire dynamics. Everything that follows is either
a consequence of it or a statement about how to apply it.

\section{Properties of the operator}\label{sec:props}

\begin{theorem}[Exact conservation]\label{thm:cons}
For all $n \in \N$ and $k \geq 1$, $\sum_{i=1}^{k} R_i(n) = n$.
\end{theorem}
\begin{proof}
Exactly $r$ components equal $q+1$ and $k-r$ equal $q$, so the sum is
$r(q+1) + (k-r)q = kq + r = n$.
\end{proof}

\begin{theorem}[Bounded imbalance]\label{thm:bal}
For all $n \in \N$ and $k \geq 1$,
$\max_i R_i(n) - \min_i R_i(n) \leq 1$. Moreover, no rule satisfying
Theorem~\ref{thm:cons} with integer components attains a strictly smaller
maximum imbalance.
\end{theorem}
\begin{proof}
Every component is $q$ or $q+1$ by \eqref{eq:rule}, so the difference is at
most one. For the second claim: if all $k$ components were equal to some $m$
then $n = km$, which fails for $k \nmid n$; and if the maximum were $q$ then
$\sum_i R_i \leq kq < n$, contradicting Theorem~\ref{thm:cons}. Hence some
component is at least $q+1$ and some at most $q$ whenever $r > 0$, and the
imbalance of one is forced.
\end{proof}

\begin{theorem}[Determinism]
For a fixed ordered neighbourhood, $\R(n)$ is unique for every $n$.
\end{theorem}
\begin{proof}
Euclidean division determines $(q,r)$ uniquely; the fixed ordering then
determines which positions receive the extra unit.
\end{proof}

\begin{proposition}[Component bounds]
For all $i$, $\lfloor n/k \rfloor \leq R_i(n) \leq \lceil n/k \rceil$.
\end{proposition}
\begin{proof}
Immediate from \eqref{eq:rule}: components are $q$ or $q+1$, and $q+1 =
\lceil n/k \rceil$ when $r>0$.
\end{proof}

\begin{proposition}[Translation invariance]\label{prop:shift}
For all $n \in \N$ and $k \geq 1$, $\R(n+k) = \R(n) + \mathbf{1}$, where
$\mathbf{1}$ is the all-ones vector.
\end{proposition}
\begin{proof}
If $n = kq + r$ with $0 \leq r < k$ then $n + k = k(q+1) + r$, and the same
$r$ positions receive the extra unit in both cases.
\end{proof}

\begin{theorem}[Local state conservation]\label{thm:local}
Let $S = \Om_c + \sum_{i=1}^{k} \Om_i$ be the total over the centre and its
neighbours, counting each coordinate once. Under \eqref{eq:transition},
$S' = S$.
\end{theorem}
\begin{proof}
$S' = (\Om_c - n) + \sum_i (\Om_i + R_i(n)) = S - n + \sum_i R_i(n) = S$ by
Theorem~\ref{thm:cons}.
\end{proof}

\begin{proposition}[Source non-negativity]
If $0 \leq n \leq \Om_c$ then $\Om_c' \geq 0$.
\end{proposition}
\begin{proof}
$\Om_c' = \Om_c - n \geq 0$ by hypothesis.
\end{proof}

\begin{remark}[Admissibility is a precondition, not a conclusion]
The transition is defined only for $n \leq \Om_c$. A request exceeding the
source must be refused without partial mutation; the reference implementation
raises before any write occurs (\S\ref{sec:impl}).
\end{remark}

\section{Conservation over a lattice}\label{sec:lattice}

Theorem~\ref{thm:local} concerns one neighbourhood. Applying the transition at
every cell of a lattice requires a decision about update order, and the two
natural choices give different dynamics. Both are conservative.

\subsection{Sequential update}

Let $V$ be finite and enumerate it as $v^{(1)}, \ldots, v^{(m)}$. A sequential
sweep applies the transition at $v^{(1)}$, then $v^{(2)}$, and so on, each
using the state left by its predecessors.

\begin{theorem}[Global conservation, sequential]\label{thm:seq}
Let $T$ be the total over a finite lattice $V$, with a boundary policy under
which every neighbour of every $v \in V$ is either in $V$ or is a coordinate
that the policy maps into $V$. Then one sequential sweep preserves $T$
exactly.
\end{theorem}
\begin{proof}
Each individual transition preserves the total over the whole lattice: it
decreases one cell by $n$ and increases others by amounts summing to $n$
(Theorem~\ref{thm:cons}), and under the stated policy no write leaves $V$.
A sweep is a finite composition of total-preserving maps, hence preserves $T$.
\end{proof}

\subsection{Simultaneous update}

A simultaneous sweep computes every transfer from a frozen snapshot and
applies all of them at once, which is what a parallel machine does.

\begin{theorem}[Global conservation, simultaneous]\label{thm:sim}
Let $\Om$ be a finitely supported state, let $a(v)$ satisfy $0 \leq a(v) \leq
\Om(v)$ for all $v$, and let one simultaneous sweep route $a(v)$ from each $v$
to its axial neighbours. Then the total is unchanged.
\end{theorem}
\begin{proof}
For each $v$ the outgoing shares sum to $a(v)$ by Theorem~\ref{thm:cons} and
the retained quantity is $\Om(v) - a(v)$; the contribution of $v$ to the new
total is therefore $a(v) + (\Om(v) - a(v)) = \Om(v)$. Summing over $v$ gives
the result.
\end{proof}

\begin{proposition}[Order independence]\label{prop:order}
The state produced by a simultaneous sweep does not depend on the order in
which cells are visited.
\end{proposition}
\begin{proof}
Every cell reads the pre-sweep snapshot and writes into a fresh accumulator, so
the map from old state to new state is a sum of independent per-cell
contributions.
\end{proof}

Proposition~\ref{prop:order} is the formal content of the claim that the
simultaneous variant is the one a parallel implementation would realise. The
sequential variant is not order independent, and the reference implementation
exposes both so that the difference is visible rather than assumed.

\section{Boundaries and coincident neighbours}\label{sec:boundary}

\subsection{Boundary policies}

A finite domain requires an explicit policy; without one, routing at the edge
writes to coordinates that no level owns, and the state leaks away silently
(we exhibit this in \S\ref{sec:impl}). We distinguish three.

\begin{definition}
A \emph{periodic} policy identifies coordinates modulo extents
$(L_1,\ldots,L_d)$ and maps every neighbour into the domain. An
\emph{absorbing} policy treats coordinates outside the extents as absent, so a
routing step that would write outside is refused. An \emph{unbounded} policy
places no restriction and lets the support grow.
\end{definition}

Under the periodic policy Theorem~\ref{thm:seq} applies directly. Under the
absorbing policy a step at a face cell is undefined and must be refused; the
reference implementation raises rather than dropping units, because dropping
would violate conservation silently.

\subsection{Coincident neighbours}

On a domain small enough, two neighbours of the same cell can be the same
coordinate. For $d = 4$ with extents $(2,2,2,2)$ and $s = 1$, the displacement
$+e_i$ and $-e_i$ both land on the same cell, so the eight-entry neighbourhood
contains four distinct coordinates, each appearing twice.

\begin{lemma}[Accumulation]\label{lem:coincide}
Let $N_s(v)$ contain a coordinate $u$ with multiplicity $m \geq 1$. Under
\eqref{eq:transition} the cell $u$ receives $\sum_{j : N_j(v) = u} R_j(n)$,
and the sum of all receipts over distinct destinations is still $n$.
\end{lemma}
\begin{proof}
The transition adds $R_j(n)$ to the coordinate $N_j(v)$ for each $j$; if
several $j$ name the same $u$ the additions accumulate. The total added over
all $j$ is $n$ by Theorem~\ref{thm:cons}, and grouping by distinct destination
does not change a sum.
\end{proof}

\begin{corollary}[Invariant survives coincidence]\label{cor:coincide}
Let $S$ be the total over the \emph{set} $\{v\} \cup \{N_j(v)\}$, counting each
distinct coordinate once. Then $S' = S$.
\end{corollary}
\begin{proof}
Lemma~\ref{lem:coincide} gives the receipts; the source loses $n$; the
distinct-destination total gains $n$. The argument of
Theorem~\ref{thm:local} applies with the sum taken over distinct coordinates.
\end{proof}

Corollary~\ref{cor:coincide} is a small result with a practical consequence: a
conservation check that sums a raw neighbour list is incorrect on any domain
where neighbours can coincide, and it fails in the direction that matters ---
it reports a violation where there is none. The reference implementation
de-duplicates the scope and accumulates shares, which is what makes the check
trustworthy on small domains.

\section{Multilevel transfer}\label{sec:multi}

A hierarchy groups $m$ fine cells under one coarse cell. Two operators move
state between levels, and both are exact on integers.

\begin{definition}[Hierarchy]
Let the fine domain have extents $L_j$ divisible by block factors
$b_j \geq 1$, let $m = \prod_j b_j$, and let the coarse domain have extents
$L_j / b_j$. The children of a coarse cell $p$ are the $m$ fine cells
$\{\,p \circ b + t : 0 \leq t_j < b_j\,\}$.
\end{definition}

\begin{definition}[Restriction and prolongation]
Restriction sets $\Om_{\mathrm{coarse}}(p) = \sum_{c \in \mathrm{children}(p)}
\Om(c)$. Prolongation sets the children of $p$ to $\R(\Om_{\mathrm{coarse}}(p))$
with $k = m$.
\end{definition}

\begin{proposition}[Restriction is exact]
Restriction preserves the total.
\end{proposition}
\begin{proof}
The children of distinct parents are disjoint and their union is the fine
domain, so the sum of parent values equals the sum of child values.
\end{proof}

\begin{proposition}[Prolongation is exact]
Prolongation preserves the total.
\end{proposition}
\begin{proof}
Each parent's value is distributed among exactly $m$ children by the operator
with $k = m$, so the children of $p$ sum to $\Om_{\mathrm{coarse}}(p)$
(Theorem~\ref{thm:cons}).
\end{proof}

\begin{theorem}[V-cycle conservation]\label{thm:vcycle}
A cycle consisting of restriction, any number of simultaneous sweeps at the
coarse level, and prolongation preserves the fine total exactly.
\end{theorem}
\begin{proof}
Restriction preserves the total; each coarse sweep preserves the coarse total
by Theorem~\ref{thm:sim}; prolongation preserves the total. Composing three
total-preserving maps preserves the total.
\end{proof}

Theorem~\ref{thm:vcycle} is the reason the multilevel extension is worth having.
A multilevel scheme is often where exactness is quietly lost, because the
transfer operators are chosen for smoothness rather than for conservation. Here
both transfers are exact by construction, so the guarantee is inherited rather
than arranged.

\begin{remark}[A boundary policy is not optional]
Theorem~\ref{thm:vcycle} requires the coarse level to have a boundary policy.
If the coarse routing is performed on an unbounded domain, mass is written to
coordinates outside the coarse domain, no level owns them, and the total falls.
The reference implementation refuses to construct a hierarchy over an unbounded
lattice for this reason, and its cycle re-checks the total rather than
trusting the argument.
\end{remark}

\section{Reference implementation and verification}\label{sec:impl}

\subsection{Structure}

The implementation is a Python package of seven modules with a strictly
downward dependency order: lattice memory, then the routing engine, then the
multilevel layer, with an assembler and machine on top.

\emph{Lattice memory} holds a sparse map from coordinate to non-negative
integer and resolves every coordinate through a declared boundary policy.
Totals de-duplicate their argument, so a caller cannot accidentally count a
cell twice.

\emph{The engine} implements the three operations of the execution model:
CONNECT binds an ordered neighbourhood and records the invariant baseline,
RESONATE performs the transition of \eqref{eq:transition}, and COLLAPSE
re-computes the scope total and compares it with the baseline. COLLAPSE
mutates nothing. Both update semantics of \S\ref{sec:lattice} are provided.

\emph{The hierarchy} implements restriction, prolongation and the V-cycle of
\S\ref{sec:multi}, and re-checks the total at the end of each cycle.

\emph{The machine} executes a 17-mnemonic integer instruction set with an
assembler that range-checks every operand against its declared width. COLLAPSE
is a trapping instruction: on a failed invariant the machine halts and records
the violation, so the check lives in the instruction stream rather than in the
test harness.

\subsection{Verification}

The suite comprises 88 tests and runs in about six seconds on the machine used
for this paper. It is included in the repository and reproduces every number
below with a single command.

\begin{table}[h]
\centering
\small
\begin{tabular}{@{}p{5.4cm}p{5.6cm}p{3.2cm}@{}}
\toprule
Property & Scope of the check & Result \\
\midrule
Operator conservation & exhaustive, $k = 2..10$, $n = 0..5000$ (45{,}009 cases) & exact \\
Bounded imbalance & exhaustive, same range & $\leq 1$ \\
Component bounds & exhaustive, same range & holds \\
Remainder count & exhaustive, same range & holds \\
Translation invariance & $n \in \{0,3,27,1000\}$, $k = 8$ & holds \\
\midrule
Local conservation & 5{,}000 randomised routes, $n \leq 10^6$ & 0 failures \\
Global conservation, sequential & tori of $4^4$, $8^4$, $10^4$ cells & exact \\
Global conservation, simultaneous & tori of $4^4$, $8^4$, $10^4$ cells & exact \\
Order independence & forwards vs.\ reversed visit order & identical state \\
Spacing invariance & $s \in \{1,2,3\}$ on an $8^4$ torus & exact \\
\midrule
Coincident neighbours & $2^4$ torus, 8 links, 4 distinct cells & invariant holds \\
MMU aliasing & two coordinates to one frame & rejected \\
\midrule
V-cycle conservation & extents 4 and 8, repeated cycles & exact \\
V-cycle without boundary & coarse level unbounded & refused, loss reported \\
\midrule
Assembler and machine & 22 cases incl.\ a trapping program & all pass \\
\bottomrule
\end{tabular}
\caption{Verification scope. Every row is a test in the repository; the test
log is generated by the runner rather than written by hand.}
\label{tab:verify}
\end{table}

\subsection{Measured behaviour}

Table~\ref{tab:timing} reports timings of the Python reference implementation.
These figures characterise this implementation and nothing else. They are not a
comparison with any other system, and they carry no information about what a
native implementation or a dedicated circuit would achieve.

\begin{table}[h]
\centering
\small
\begin{tabular}{@{}lrr@{}}
\toprule
Operation & Cells & Rate \\
\midrule
CONNECT + RESONATE pairs, unbounded lattice & --- & 9{,}700 routes/s \\
Sequential sweep, periodic $4^4$ torus & 256 & 4{,}917 cells/s \\
Sequential sweep, periodic $8^4$ torus & 4{,}096 & 4{,}493 cells/s \\
Sequential sweep, periodic $10^4$ torus & 10{,}000 & 3{,}868 cells/s \\
Simultaneous sweep, periodic $8^4$ torus & 4{,}096 & 10{,}175 cells/s \\
Simultaneous sweep, periodic $10^4$ torus & 10{,}000 & 9{,}126 cells/s \\
V-cycle, block $2^4$, extent 4 & 256 & 3.9 ms/cycle \\
V-cycle, block $2^4$, extent 8 & 4{,}096 & 72.6 ms/cycle \\
\bottomrule
\end{tabular}
\caption{Timings of the Python reference implementation, single core,
measured by \texttt{examples/bench.py}. The simultaneous sweep is faster
because it writes one accumulator per target rather than re-reading cells.
All runs conserved the total exactly.}
\label{tab:timing}
\end{table}

\subsection{A leak that the tests catch}

The multilevel module is instructive about why the boundary policy must be
explicit. Constructing a hierarchy over a periodic $4^4$ domain, then replacing
the coarse memory with an unbounded one and running a cycle, halves the total:
$2{,}560 \to 1{,}280$. The coarse routing writes to parent coordinates outside
the coarse domain; the unbounded policy creates them; no level sums them; the
state is gone, with no error raised by the arithmetic. The cycle compares the
fine total before and after and raises, which converts a silent corruption into
a reported failure. This is the only reason the module re-checks a total that
Theorem~\ref{thm:vcycle} already guarantees: the theorem's hypotheses include a
condition that a caller can violate by configuration.

\section{Related work}\label{sec:related}

\paragraph{Apportionment.} The rule \eqref{eq:rule} is the largest-remainder
method: allocate the floor to each party, then distribute the remaining seats
to those with the largest fractional parts. Its properties, its history and its
pathologies are treated by Balinski and Young \cite{balinski1982}. Their
analysis is the correct reference point for the per-cell results of
\S\ref{sec:props}, and we do not claim those results as new. What differs here
is the setting: the rule is used as a local transport operator on a lattice
rather than as a one-shot allocation among a fixed set of parties, which is
what makes the lattice-level and multilevel results of
\S\ref{sec:lattice}--\ref{sec:multi} available.

\paragraph{Lattice gas and cellular automata.} Conservative lattice automata
with integer states have a long history, including lattice gas automata for
fluid simulation \cite{frisch1986} and the broader study of cellular automata
\cite{wolfram2002}. Those models typically conserve a vector of discrete
quantities through collision rules that are local and reversible. The present
operator is simpler --- it conserves a single scalar and is not reversible ---
and should be read as a transport primitive rather than as a physical model.

\paragraph{Integer and exact arithmetic.} The absence of rounding error is a
property of integer arithmetic and not an achievement of this construction. The
relevant comparison is with fixed-point and rational arithmetic, where exactness
is also available and where the cost of wider representations has been studied
extensively \cite{yates2013}.

\paragraph{Multigrid.} Restriction and prolongation are the standard multigrid
transfer operators \cite{briggs2000}, usually chosen as weighted averages and
adjoints. Choosing them to be exact on integers instead --- a sum and a split
--- gives up smoothing quality in exchange for exact conservation. We make no
claim that the resulting scheme is an efficient solver; we claim only that its
cycles are conservative, which is what Theorem~\ref{thm:vcycle} establishes.

\section{Limitations}\label{sec:limits}

\begin{enumerate}[leftmargin=1.4em]
\item \textbf{No performance result.} There is no comparison with any existing
system, and the timings in Table~\ref{tab:timing} describe a Python reference
implementation only.

\item \textbf{No hardware.} The instruction set of \S\ref{sec:impl} is a
specification. No circuit has been built, synthesised or measured, and the
paper draws no conclusion about feasibility, area, frequency or power.

\item \textbf{No physical interpretation.} The lattice is a data structure.
Nothing here establishes that it models spacetime, energy, thermodynamics or
entropy, and no physical quantity is claimed to be conserved.

\item \textbf{The operator is a smoother, not a solver.} Repeated sweeps drive
a periodic domain towards the uniform state. That is the behaviour of a
conservative averaging operator and it does not solve any physical problem.

\item \textbf{Novelty is limited and stated.} The routing rule is the
largest-remainder rule. The per-cell properties of \S\ref{sec:props} follow
from it directly. The contributions this paper can defend are the lattice-level
results, the coincident-neighbour lemma, the multilevel conservation theorem,
and the verified implementation.

\item \textbf{Verification is testing, not proof.} The exhaustive checks cover
$k \leq 10$ and $n \leq 5000$; the randomised checks cover larger values. The
theorems are proved, but the correspondence between the theorems and the code is
established by test, not by formal refinement.

\item \textbf{Sequential update is order dependent.} Proposition~\ref{prop:order}
holds only for the simultaneous variant. Any result that depends on the
sequential dynamics must state the visiting order.

\item \textbf{Absorbing boundaries refuse rather than redistribute.} A routing
step at a face cell is an error under the absorbing policy. Whether that is the
right semantics for a bounded problem is a modelling question this paper does
not settle.
\end{enumerate}

\section{Future work}\label{sec:future}

One item is a genuine research question and the rest are engineering.

\paragraph{Classification of the rule family.} Characterise all maps
$\R : \N \to \N^k$ that are conservative (Theorem~\ref{thm:cons}), balanced to
within one (Theorem~\ref{thm:bal}), deterministic, and translation invariant in
the sense of Proposition~\ref{prop:shift}. The conjecture is that the family is
exactly the orderings of \eqref{eq:rule}: that any rule with these four
properties must be the largest-remainder rule composed with a permutation of
the destinations. If that holds, the operator is not one choice among many but
the only one, and the paper's contribution becomes a characterisation rather
than a description. This is the single result most likely to make the work
publishable in a discrete-mathematics venue, and it is the first item in the
project's plan.

\paragraph{Convergence of repeated sweeps.} The measured spread decreases
monotonically under repeated simultaneous sweeps. The rate, the conditions for
reaching uniformity exactly rather than asymptotically, and the dependence on
the spacing $s$ are open.

\paragraph{Boundary theory.} Reflecting and absorbing boundaries interact with
the invariant in ways not analysed here; the absorbing case currently refuses
rather than redistributes.

\paragraph{A native implementation with a fair baseline.} The performance
question cannot be answered in Python. It requires a C or Rust kernel, a
cache-blocked integer baseline solving the same problem, and measurement with
hardware counters rather than synthetic ones. The project plan treats this as a
gate: if the integer formulation does not beat the baseline on a real workload,
the accelerator line of work ends and the mathematics is published on its own.

\paragraph{A hardware prototype, conditionally.} The arithmetic is
gate-friendly: the routing is a three-bit shift, a one-hot remainder mask and
$k$ adders, and the invariant check is an adder tree and a comparator. Whether
that is worth building depends entirely on the previous item.

\section{Conclusion}\label{sec:conclusion}

The operator studied here is elementary: split an integer by the neighbourhood
size, distribute the quotient, hand the remainder out one unit at a time. Its
per-cell properties are immediate consequences of Euclid's division lemma, and
the rule itself is two centuries old.

What is not immediate is how the operator behaves when it is the local rule of a
whole lattice. Conservation survives both update orders; it survives
coincident neighbours on a small periodic domain, provided the conservation
scope is de-duplicated; and it survives a two-level restriction--prolongation
cycle, because both transfer operators are exact rather than approximate. Those
results, together with a verified implementation and an explicit statement of
what is not claimed, are what this paper offers.

The honest summary is that the mathematics is small and sound, and the
engineering question --- whether conservative integer transport is worth
accelerating --- is open and answerable. Answering it is the next step, and it
is the step that decides whether the work matters beyond its own correctness.

\appendix

\section{Withdrawn claims}\label{app:withdrawn}

Earlier versions of this work, circulated as a thesis, a whitepaper and a
manuscript submission, contained claims that the accompanying code did not
support. They are recorded here so that anyone comparing versions understands
what changed. Each was checked against the implementation and withdrawn.

\begin{enumerate}[leftmargin=1.4em]
\item \textbf{Reductions of 87.50\%, 89.47\% and 94.44\% in data-bus traffic.}
These were ratios of transaction counters incremented by hand-written
constants in the measurement code, not measurements of memory traffic. The
three figures came from three different constant assignments and none was
reproducible from the shipped code. Withdrawn in full. No replacement figure is
offered, because none was measured.

\item \textbf{``0.00\% rounding drift'' as an architectural property.} True of
any integer arithmetic and not specific to this construction. Restated as a
property of the representation rather than a result.

\item \textbf{Axial routing described as four-dimensional while the code used
Moore neighbourhoods.} The simulations returned four
orthogonal plus four diagonal cells. The code and the description disagreed;
the description is corrected here and the neighbourhood is now derived from the
dimension.

\item \textbf{A self-hosting bare-metal compiler and kernel.} The artifacts were
a pseudo-assembly listing with no executor, an assembler that compiled a
hardcoded three-line string and could not parse that listing, and a 37-byte
binary. Withdrawn. What replaces it is a specification and a working assembler
and machine, with no claim of hardware realisation.

\item \textbf{Physical interpretation: imaginary-mass modes, a cosmological
singularity reset, and a consciousness decomposition of an energy equation.}
These have no counterpart in the code and no derivation. Withdrawn in full.

\item \textbf{Applications to planetary atmospheric modelling and exascale
physics.} The demonstrated operator is a conservative smoother. No atmospheric
model was built or validated. Withdrawn.

\item \textbf{``Neuromorphic'' and ``unified field'' as descriptors.} The
construction involves no spiking dynamics, no plasticity and no field theory.
The names are retained only where they refer to the historical project and
carry no technical claim.
\end{enumerate}

\begin{thebibliography}{9}

\bibitem{balinski1982}
M.~L. Balinski and H.~P. Young,
\emph{Fair Representation: Meeting the Ideal of One Man, One Vote},
Yale University Press, 1982.

\bibitem{frisch1986}
U.~Frisch, B.~Hasslacher and Y.~Pomeau,
``Lattice-gas automata for the Navier--Stokes equation,''
\emph{Physical Review Letters}, vol.~56, no.~14, pp.~1505--1508, 1986.

\bibitem{wolfram2002}
S.~Wolfram,
\emph{A New Kind of Science}, Wolfram Media, 2002.

\bibitem{briggs2000}
W.~L. Briggs, V.~E. Henson and S.~F. McCormick,
\emph{A Multigrid Tutorial}, 2nd ed., SIAM, 2000.

\bibitem{yates2013}
R.~Yates,
``Fixed-point arithmetic: an introduction,''
Digital Signal Labs, technical reference, 2013.

\bibitem{knuth1997}
D.~E. Knuth,
\emph{The Art of Computer Programming, Volume 1: Fundamental Algorithms},
3rd ed., Addison--Wesley, 1997.

\end{thebibliography}

\end{document}
